\subsubsection{Quantum Advantage and Importance}\label{ex:quantum-advantage-and-importance}

Deutsch's algorithm is the first example of a quantum algorithm that provably outperforms every deterministic classical algorithm for the same problem. Although the speedup is modest, it introduces several of the fundamental ideas that appear in many later quantum algorithms.

\highspace
\begin{flushleft}
    \textcolor{Green3}{\faIcon{balance-scale} \textbf{One Quantum Query vs Two Classical Queries}}
\end{flushleft}
The goal is to determine whether an unknown Boolean function:
\begin{equation*}
    f:\{0,1\}\to\{0,1\}
\end{equation*}
Is \textbf{constant} or \textbf{balanced}.

\highspace
\textcolor{Red2}{\faIcon{times-circle}} A deterministic classical algorithm must evaluate $f(0)$ and $f(1)$ (\textbf{two queries}) to determine whether $f$ is constant or balanced. Because knowing only one output is never sufficient.

\highspace
\textcolor{Green3}{\faIcon{check-circle}} Deutsch's algorithm instead prepares a superposition of both inputs:
\begin{equation*}
    \frac{\ket{0} + \ket{1}}{\sqrt{2}}
\end{equation*}
Queries the oracle \textbf{once}, and uses phase kickback together with quantum interference to determine whether:
\begin{equation*}
    f(0) = f(1) \quad \text{(constant)}
\end{equation*}
Or:
\begin{equation*}
    f(0) \neq f(1) \quad \text{(balanced)}
\end{equation*}
Thus, Deutsch's algorithm achieves a \textbf{quantum advantage} over classical algorithms by requiring only one query to the oracle instead of two.

\highspace
\begin{flushleft}
    \textcolor{Green3}{\faIcon{question-circle} \textbf{Why This Is a Query-Complexity Advantage}}
\end{flushleft}
Notice that the algorithm does \textbf{not} reduce the number of elementary quantum gates. Preparing the state, applying Hadamard gates, and performing the measurement all require additional operations.

\highspace
Instead, the improvement concerns the number of times the algorithm must access the \textbf{oracle} $U_f$. This is called the \textbf{query complexity} (or \textbf{oracle complexity}) of the problem. Many quantum algorithms are analyzed in exactly this way: ordinary gates are considered inexpensive; oracle evaluations represent the costly resource. Deutsch's algorithm demonstrates that quantum mechanics can reduce the number of oracle queries by exploiting superposition, phase kickback, and interference.

\newpage

\begin{flushleft}
    \textcolor{Green3}{\faIcon{question-circle} \textbf{Why the Algorithm Is Historically Important}}
\end{flushleft}
The computational speedup itself is small:
\begin{equation*}
    2 \text{ classical queries} \quad \text{vs} \quad 1 \text{ quantum query}
\end{equation*}
Nevertheless, Deutsch's algorithm is historically significant because it is the \hl{first algorithm to demonstrate a genuine quantum computational advantage.} More importantly, it introduces a design pattern that reappears throughout quantum computing:
\begin{enumerate}
    \item Prepare a superposition of many inputs;
    \item Encode information into relative phases using an oracle;
    \item Use interference to amplify the desired information;
    \item Measure only a global property of the function.
\end{enumerate}
Instead of learning every function value individually, the \hl{algorithm extracts only the property that is needed.}

\highspace
\begin{flushleft}
    \textcolor{Green3}{\faIcon{link} \textbf{Connection to Later Quantum Algorithms}}
\end{flushleft}
Many of the best-known quantum algorithms follow exactly the same principles introduced by Deutsch's algorithm.
\begin{itemize}
    \item \important{Deutsch-Jozsa Algorithm}. The Deutsch-Jozsa algorithm \textbf{generalizes the problem} from one input bit to $n$ input bits:
    \begin{equation*}
        f : \{0,1\}^n \to \{0,1\}
    \end{equation*}
    While maintaining the promise that the function is either constant or balanced. Remarkably, the quantum algorithm still requires only \textbf{one oracle query}, whereas a deterministic classical algorithm may require $2^{n-1}+1$ queries in the worst case.


    \item \important{Bernstein-Vazirani Algorithm}. The Bernstein-Vazirani algorithm also relies on phase kickback and interference. Instead of determining whether a function is constant or balanced, its goal is to \textbf{determine an unknown $n$-bit string hidden inside the oracle}. A classical deterministic algorithm requires $n$ oracle queries, while the quantum algorithm determines the entire hidden string using only \textbf{one oracle query}.
\end{itemize}